% ============================================================
% РАЗДЕЛ «ДИФФЕРЕНЦИАЛЬНАЯ ЗАЩИТА ЛИНИИ»
% ============================================================

{\color{unimain}\subsection{Дифференциальная защита линии}}

% Макрос для вставки рисунка с подписью и индивидуальным масштабом
% Использование: \dzlfig[опции]{путь}{подпись}{метка}
% Опции по умолчанию: width=0.9\textwidth, height=0.9\textheight, keepaspectratio
\newcommand{\dzlfig}[4][]{%
  \begin{figure}[H]
    \centering
    \includegraphics[#1]{#2}
    \captionof{figure}{#3}\label{#4}
  \end{figure}%
  \medskip
}

\begin{enumerate}

% ============================================================
% ОБЩИЕ СВЕДЕНИЯ
% ============================================================
\item Общие сведения
\begin{enumerate}
\item Дифференциальная защита линии (ДЗЛ) является быстродействующей защитой с абсолютной селективностью и выполняет функцию основной токовой защиты линии от всех видов КЗ.
\item ДЗЛ предназначена для защиты двухконцевой или трехконцевой линии с/без ответвлений от всех видов КЗ и действует на отключение линии со всех сторон.
\item Принцип действия ДЗЛ основан на пофазном сравнении векторов тока по концам защищаемой линии путем расчета дифференциального тока. При выявлении дифференциального тока, превышающего заданный уставкой порог, защита формирует команду на отключение выключателей со всех сторон линии.
\item ДЗЛ состоит из полукомплектов (п/к), устанавливаемых на концах защищаемой линии и соединяемых между собой независимыми дублирующими цифровыми каналами связи (ЦКС) для обмена дискретной и аналоговой информацией.
\item ДЗЛ селективно срабатывает:
  \begin{itemize}
    \item при всех видах внутренних КЗ, в том числе при постановке линии под напряжение;
    \item при переходе внешнего КЗ во внутреннее, в том числе при реверсе мощности.
  \end{itemize}
\item Обеспечивается несрабатывание защиты:
  \begin{itemize}
    \item при постановке линии под напряжение и включении в транзит при отсутствии КЗ;
    \item при внешних КЗ, в том числе при реверсе мощности;
    \item при синхронных качаниях и в режиме асинхронного хода;
    \item при КЗ за трансформаторами ответвлений (отпаек) на линиях с ответвлениями;
    \item при бросках тока намагничивания трансформаторов ответвлений (отпаек) без КЗ;
    \item при несимметрии токов, вызванной тяговой нагрузкой;
    \item при внешних КЗ с насыщением ТТ (при соблюдении заявленных производителем требований к ТТ).
  \end{itemize}
\item ДЗЛ обеспечивает выполнение следующих требований к времени срабатывания:
  \begin{itemize}
    \item время срабатывания устройства ДЗЛ на отключение (без учета задержки в канале связи и с учетом времени работы выходных реле) при протекании двухкратного тока срабатывания – не более 45 мс;
    \item собственное время срабатывания измерительных органов ДЗЛ (без учета задержки в канале связи) при переходе внешнего КЗ во внутреннее в условиях наличия насыщения ТТ – не более 60 мс;
    \item собственное время срабатывания измерительных органов ДЗЛ при подаче двукратного тока срабатывания – не более 25 мс.
  \end{itemize}
\item ДЗЛ обеспечивает время достоверного измерения значения тока в переходных режимах, сопровождающихся насыщением ТТ не более 6 мс.
\item Полукомплекты по концам защищаемой линии обмениваются между собой выборками мгновенных значений первичных токов плеч в относительных единицах (под плечом подразумевается участок электрической цепи, соответствующий одному из выключателей присоединения). Для этого используется приведение измеренных токов к первичному базисному току $ I_{\textit{баз}} $ с учетом инвертирования токов по формуле:
  \begin{align}
    i_{\textit{*перв}}=\frac{i_{\textit{втор}} \cdot I_{\textit{ном перв}} \cdot k_{\textit{инв}}}{I_{\textit{ном втор}} \cdot I_{\textit{баз}}},
  \end{align}
  \begin{eqexpl}[7mm]
    \item{$ i_{\textit{*перв}} $} мгновенное значение первичного тока, приведенное к $ I_{\textit{баз}} $, о.е.;
    \item{$ i_{\textit{втор}} $} мгновенное значение вторичного тока, А;
    \item{$ I_{\textit{ном перв}} $} \textls[-5]{номинальный первичный ток ТТ, который задается уставками \textbf{«Iном перв В1»} и \textbf{«Iном перв В2»}, А;}
    \item{$ I_{\textit{ном втор}} $} \textls[-5]{номинальный вторичный ток ТТ, который задается уставками \textbf{«Iном втор В1»} и \textbf{«Iном втор В2»}, А;}
    \item{$ k_{\textit{инв}} $} коэффициент программного инвертирования тока, который равен «1» или «-1» в зависимости от выбора уставок \textbf{«Инверсия ТТ В1»} и \textbf{«Инверсия ТТ В2»};
    \item{$ I_{\textit{баз}} $} первичный базисный ток, величина которого задается уставкой \textbf{«Iбаз ДЗЛ»}, А.
  \end{eqexpl}
\item В качестве базисного тока выбирается максимальный номинальный первичный ток ТТ, установленного в одном из плеч п/к ДЗЛ. Важно отметить, что базисный ток для двухконцевой ДЗЛ должен задаваться одинаковым для обоих п/к. Для трехконцевой ДЗЛ базисный ток должен задаваться одинаковым соответственно для всех трех п/к. Также для ограничения тока небаланса в дифференциальном токе первичные номинальные токи ТТ не должны отличаться более чем в 2,5 раза как по концам линии, так и между ТТ двух выключателей линии в пределах одной подстанции. Параметры настройки ТТ ДЗЛ приведены в таблице \ref{desc_dzl:tbl1}.
  \begingroup
  \small
  \setlength{\extrarowheight}{0.06cm}
  \begin{longtable}{|>{\centering\arraybackslash}m{6.4cm}|>{\centering\arraybackslash}m{4cm}|>{\centering\arraybackslash}m{2.8cm}|>{\centering\arraybackslash}m{2.8cm}|}
    \caption{Перечень параметров настройки ТТ ДЗЛ\hfill\vspace{-0.5\baselineskip}}\label{desc_dzl:tbl1}\\
    \hline
    \rowcolor{gray!30} Наименование уставки & Диапазон & Шаг & Значение по умолчанию \\
    \hline
    \endfirsthead
    \caption*{\hspace{3pt}\emph{Продолжение таблицы \ref{desc_dzl:tbl1}}\hfill\ }\vspace{-0.5\baselineskip} \\
    \hline
    \rowcolor{gray!30} Наименование уставки & Диапазон & Шаг & Значение по умолчанию \\
    \hline
    \endhead
    \centering\arraybackslash Iном перв В1  & \centering\arraybackslash 0…1000000 А & \centering\arraybackslash 1 А & \centering\arraybackslash 1000 А \\
    \hline
    \centering\arraybackslash Iном перв В2  & \centering\arraybackslash 0…1000000 А & \centering\arraybackslash 1 А & \centering\arraybackslash 1000 А \\
    \hline
    \centering\arraybackslash Iном втор В1  & \centering\arraybackslash 1…5 А & \centering\arraybackslash 1 А & \centering\arraybackslash 5 А \\
    \hline
    \centering\arraybackslash Iном втор В2  & \centering\arraybackslash 1…5 А & \centering\arraybackslash 1 А & \centering\arraybackslash 5 А \\
    \hline
    \centering\arraybackslash Инверсия ТТ В1  & \centering\arraybackslash Введено/Выведено & \centering\arraybackslash - & \centering\arraybackslash Выведено \\
    \hline
    \centering\arraybackslash Инверсия ТТ В2  & \centering\arraybackslash Введено/Выведено & \centering\arraybackslash - & \centering\arraybackslash Выведено \\
    \hline
  \end{longtable}
  \endgroup
\item В рамках каждого функционального блока реализован возможность его вывода с помощью сигналов \textit{«Вывод ДЗЛ»} и \textit{«Вывод терминала»} или при установке соответствующего положения программного ключа \textbf{«ДЗЛ»}.
\end{enumerate}

% ============================================================
% РЕЖИМЫ РАБОТЫ ДЗЛ
% ============================================================
\item Режимы работы ДЗЛ
\begin{enumerate}
\item Каждый п/к ДЗЛ имеет следующие режимы работы:
  \begin{itemize}
    \item «Сопряжение»;
    \item «Работа»;
    \item «Тест».
  \end{itemize}
\item Режим «Сопряжение» является временным служебным режимом, необходимым для сопряжения п/к ДЗЛ. Инициализация сопряжения осуществляется из соответствующего меню программного обеспечения по работе с устройством М500. В данном режиме производится автоматическая блокировка ДЗЛ с обнулением значений принимаемых токов с удаленного п/к ДЗЛ. В рамках данного режима каждый п/к ДЗЛ обменивается идентификатором «ID\_Tx», который определяется уставкой \textbf{«ID Tx»}. Каждый п/к ДЗЛ проверяет соответствие установленного параметра \textbf{«ID Rx»} и принимаемого идентификатора данных «ID\_Tx» от удаленного п/к ДЗЛ.
\item В случае успешного сопряжения устанавливается сигнал \textit{"успешное сопряжение"} и следует переход в режим «Работа». В случае, если в течение 5 секунд отсутствует факт успешного сопряжения, то устанавливается сигнал \textit{"неуспешное сопряжение"} и работа ДЗЛ блокируется.
\item Режим «Работа» является основным режимом работы. П/к ДЗЛ находится в данном режиме работы при отсутствии введенного режима тестирования. В зависимости от уставки \textbf{«РежСинхИзм»} возможны следующие режимы синхронизации:
  \begin{itemize}
    \item «Внутренняя». Установлен внутренний источник синхронизации измерений (Эхо-метод);
    \item «Внешняя». Установлен внешний источник синхронизации измерений (1PPS).
  \end{itemize}
\item Режим «Тест» является вспомогательным режимом работы. Данный режим работы необходим для проведения наладочных работ при отсутствии возможности выдачи тестовых токов на оба п/к. В данном режиме работы для проверки работы функции ДЗЛ используются токи одного из п/к. Режим «Тест» имеет следующие особенности работы п/к ДЗЛ:
  \begin{itemize}
    \item режим «Тест» работает аналогично режиму «Работа» за исключением необходимости ретрансляции получаемых токов удаленным п/к (A->B, B->C, C->A) в токи соседних фаз. Например, ток фазы А должен передаваться в фазе В, B->C, C->A;
    \item при расчете дифференциального и тормозного тока в режиме «Тест» учитываются особенности передачи тестовых измерений. Например, дифференциальный ток ф.А определяется по формуле (2) с учетом того, что суммируются ток ф.А локального п/к ДЗЛ и ток ф.В удаленного п/к ДЗЛ, для диф. тока ф.В - ток ф. В локального п/к и ф.С удаленного, для диф. тока ф.С - ток ф. С локального п/к ДЗЛ и ток ф.А удаленного п/к ДЗЛ.
  \end{itemize}
\end{enumerate}

% ============================================================
% КОМПЕНСАЦИЯ ЕМКОСТНОГО ТОКА
% ============================================================
\item Компенсация емкостного тока
\begin{enumerate}
\item При применении ДЗЛ на кабельных и протяженных воздушных линиях может возникать значительный небаланс в дифференциальном токе, обусловленный влиянием емкостной составляющей. Для устранения этого влияния предусмотрена функция компенсации емкостного тока, которая вводится программным ключом \textbf{«Комп. ЕТ»}.
\item По каналам связи между п/к ДЗЛ передаются уже компенсированные фазные токи.
\item Компенсация емкостного тока осуществляется пофазно в мгновенных величинах по формулам:
  \begin{align}
    i_{\textit{a комп}} = i_{\textit{a}} - i_{\textit{a емк}}, \quad
    i_{\textit{b комп}} = i_{\textit{b}} - i_{\textit{b емк}}, \quad
    i_{\textit{c комп}} = i_{\textit{c}} - i_{\textit{c емк}}
  \end{align}
  \begin{eqexpl}[17mm]
    \item{$ i_{\textit{a}}, i_{\textit{b}}, i_{\textit{c}} $} мгновенные токи соответствующих фаз, А;
    \item{$ i_{\textit{a емк}}, i_{\textit{b емк}}, i_{\textit{c емк}} $} мгновенные емкостные токи соответствующих фаз, А.
  \end{eqexpl}
\item В каждом п/к ДЗЛ расчёт полного ёмкостного тока линии выполняется по Π-образной схеме замещения на основе суммарной реактивной проводимости линии (прямая и нулевая последовательности), измеренных значений фазных напряжений, с учетом количества выключателей присоединения и их текущего положения. Мгновенный емкостной ток фаз рассчитывается по формуле:
  \begin{align}
    \left(\begin{array}{c}i_{\textit{a емк}}\\ i_{\textit{b емк}}\\i_{\textit{c емк}}\end{array}\right)=
    \left(\begin{array}{ccc}
      \frac{2C_1+C_0}{3} & \frac{C_0-C_1}{3} & \frac{C_0-C_1}{3} \\
      \frac{C_0-C_1}{3} & \frac{2C_1+C_0}{3} & \frac{C_0-C_1}{3}\\
      \frac{C_0-C_1}{3} & \frac{C_0-C_1}{3} & \frac{2C_1+C_0}{3}
    \end{array}\right)
    \cdot \left(\begin{array}{c}\frac{\partial u_a}{\partial t}\\ \frac{\partial u_b}{\partial t} \\ \frac{\partial u_c}{\partial t} \end{array}\right)
  \end{align}
  \begin{eqexpl}[25mm]
    \setlength{\itemsep}{1.05\baselineskip}
    \item{$ C_1 = \frac{10^{-6} \cdot b_1 \cdot L}{\omega} $} полная емкость линии прямой последовательности;
    \item{$ C_0 = \frac{10^{-6} \cdot b_0 \cdot L}{\omega} $} полная емкость линии нулевой последовательности;
    \item{$ \frac{\partial u_a}{\partial t}, \frac{\partial u_b}{\partial t}, \frac{\partial u_c}{\partial t} $} производные соответствующих фазных напряжений по времени;
    \item{$ b_1 $} удельная емкостная проводимость линии прямой последовательности, мкСм/км. Значение определяется уставкой \textbf{«Удельная проводимость прямой последовательности»};
    \item{$ b_0 $} удельная емкостная проводимость линии нулевой последовательности, мкСм/км. Значение определяется уставкой \textbf{«Удельная проводимость нулевой последовательности»};
    \item{$ L $} длина линии, км. Значение определяется уставкой \textbf{«Длина защищаемой линии»};
    \item{$ \omega $} угловая частота основной гармоники тока, рад/с.
  \end{eqexpl}

\item Поясняющая схема компенсации емкостного тока на двухконцевой ЛЭП представлена на рисунке \ref{pict_dzl:KET_line}.

\dzlfig[width=0.8\textwidth,height=0.8\textheight,keepaspectratio]{\fbpath/02.01 ДЗЛ/561.0651 ДЗЛ/img/KET_line.pdf}{Компенсация емкостного тока двухконцевой линии}{pict_dzl:KET_line}

\item Предусмотрен алгоритм пропорционального распределения полного емкостного тока между п/к ДЗЛ. Пропорциональность компенсации каждым п/к определяется программным ключом \textbf{«Топология ЛЭП»}. При положении ключа \textbf{«Топология ЛЭП»} = «Двухконц.» в п/к ДЗЛ с введенной компенсацией компенсируется половина емкостного тока, при положении ключа \textbf{«Топология ЛЭП»} = «Трехконц.» соответственно треть.

\item Компенсация емкостных токов возможна только при условии правильного измерения всех фазных напряжений линии. При срабатывании БНН в одном из п/к или при отключенном положении выключателей линии в случае установки рабочего ТН на шинах, функция компенсации емкостного тока на этом конце автоматически выводится. Место установки рабочего ТН определяется положением программного переключателя \textbf{«Место ТН»}.

\item При автоматическом или оперативном выводе компенсации в одном из п/к ДЗЛ компенсация выводится во всех п/к.
\item Компенсация емкостного тока представлена на рисунке \ref{pict_dzl:UKET}.
  \dzlfig[width=0.9\textwidth,height=0.9\textheight,keepaspectratio]{\fbpath/02.01 ДЗЛ/561.0651 ДЗЛ/img/UKET.pdf}{Компенсация емкостного тока}{pict_dzl:UKET}
\item Параметры для настройки УКЕТ представлены в таблице \ref{app_set:dzl}.
\end{enumerate}

% ============================================================
% ИЗМЕРИТЕЛЬНЫЕ ОРГАНЫ ДЗЛ
% ============================================================
\item Измерительные органы ДЗЛ
\begin{enumerate}
\item Измерительные органы (ИО) ДЗЛ осуществляют обработку аналогового сигнала и сопоставление полученного результата с заданной уставкой срабатывания. На основании вычислений формируется соответствующий логический сигнал, определяющий факт срабатывания либо несрабатывания.
\item ИО ДЗЛ включает в себя следующие ИО:
  \begin{itemize}
    \item дифференциальная токовая защита с торможением (ДТЗт);
    \item дифференциальная токовая отсечка (ДТО);
    \item медленнодействующий дифференциальный орган (МДО);
    \item быстродействующий дифференциальный орган (БДО);
    \item блокировка при внешних коротких замыканиях (БВКЗ).
  \end{itemize}
\item Основными рабочими величинами ИО являются действующие значения первой гармоники дифференциального и тормозного токов.
\item Действующее значение дифференциального тока вычисляется как модуль геометрической суммы первых гармоник токов по концам защищаемой линии отдельно для каждой фазы:
  \begin{align}
    I_\textit{*Д}=\mid \dot{I}_\textit{*В1}^\textit{(А)} + \dot{I}_\textit{*B2}^\textit{(А)} + \dot{I}_\textit{*B1}^\textit{(Б)} + \dot{I}_\textit{*B2}^\textit{(Б)} + \dot{I}_\textit{*B1}^\textit{(В)} + \dot{I}_\textit{*B2}^\textit{(В)}\mid,
  \end{align}
  \begin{eqexpl}[4mm]
    \item{$ I_\textit{*Д} $} действующее значение первой гармоники дифференциального тока, о.е.;
    \item{$ \dot{I}_\textit{*В1}^\textit{(А)} $, $ \dot{I}_\textit{*В2}^\textit{(А)} $} комплексные действующие значения первой гармоники токов плеч В1 и В2 на подстанции А (локального п/к);
    \item{$ \dot{I}_\textit{*В1}^\textit{(Б)} $, $ \dot{I}_\textit{*В2}^\textit{(Б)} $} комплексные действующие значения первой гармоники токов плеч В1 и В2 на подстанции Б (удаленного п/к);
    \item{$ \dot{I}_\textit{*В1}^\textit{(В)} $, $ \dot{I}_\textit{*В2}^\textit{(В)} $} комплексные действующие значения первой гармоники токов плеч В1 и В2 на подстанции В (удаленного п/к).
  \end{eqexpl}
\item Действующее значение тормозного тока первой гармоники рассчитывается как полусумма модулей токов по концам защищаемой линии:
  \begin{align}
    I_\textit{*Т}=0,5 \cdot (\mid \dot{I}_\textit{*В1}^\textit{(А)} \mid + \mid \dot{I}_\textit{*B2}^\textit{(А)} \mid + \mid \dot{I}_\textit{*B1}^\textit{(Б)} \mid + \mid \dot{I}_\textit{*B2}^\textit{(Б)} \mid + \mid \dot{I}_\textit{*B1}^\textit{(В)} \mid + \mid \dot{I}_\textit{*B2}^\textit{(В)} \mid),
  \end{align}
  \begin{eqexpl}[4mm]
    \item{$ I_\textit{*Т} $} действующее значение первой гармоники тормозного тока, о.е.
  \end{eqexpl}
\item Измерительные органы ДЗЛ представлены на рисунке \ref{pict_dzl:IO_DZL}.
  \dzlfig[width=0.95\textwidth,height=0.95\textheight,keepaspectratio]{\fbpath/02.01 ДЗЛ/561.0651 ДЗЛ/img/IO_DZL.pdf}{Измерительные органы ДЗЛ}{pict_dzl:IO_DZL}
\item Параметры для настройки ИО ДЗЛ представлены в таблице \ref{app_set:dzl}.
\end{enumerate}

% ============================================================
% ДИФФЕРЕНЦИАЛЬНАЯ ТОКОВАЯ ЗАЩИТА С ТОРМОЖЕНИЕМ
% ============================================================
\item Дифференциальная токовая защита с торможением

\begin{enumerate}

\item ДТЗт предназначена для отключения защищаемой линии при всех видах внутренних повреждений. Ввод в работу ДТЗт осуществляется программным ключом \textbf{«ДТЗт»}.

\item Принцип работы ДТЗт основан на определении положения конца вектора дифференциального и тормозного тока относительно тормозной характеристики. При нахождении годографа тока выше тормозной характеристики формируется сигнал срабатывания ДТЗт.

\item Тормозная характеристика ДТЗт включает один горизонтальный и два наклонных участка.

\item Положение горизонтального участка определяется уставкой \textbf{«Iд нач.»}, этот участок является основным рабочим горизонтальным участком. Начало наклонных участков определяется соответствующими токами начала торможения \textbf{«Iт1»} и \textbf{«Iт2»}, а их наклон соответствующими коэффициентами торможения \textbf{«Kт1»} и \textbf{«Kт2»}.

\item Предусмотрена возможность загрубления ДТЗт при выводе УКЕТ. При загрублении вводится грубая характеристика срабатывания ДТЗт, начальный дифференциальный ток срабатывания которой определяется уставкой \textbf{«Iд нач. груб»}.

\dzlfig[width=0.65\textwidth,height=0.65\textheight,keepaspectratio]{\fbpath/02.01 ДЗЛ/561.0651 ДЗЛ/img/DTZt_torm_char.pdf}{Тормозная характеристика ДТЗт}{pict_dzl:DTZt_torm_char}

\end{enumerate}

% ============================================================
% ДИФФЕРЕНЦИАЛЬНАЯ ТОКОВАЯ ОТСЕЧКА
% ============================================================
\item Дифференциальная токовая отсечка
\begin{enumerate}
\item ДТО предназначена для быстрого отключения защищаемого объекта при внутренних повреждениях, сопровождающихся большими токами. Принцип работы ДТО основан на сравнении действующего значения первой гармоники дифференциального тока с уставкой \textbf{«Iср ДТО»}, определяющей ток срабатывания ДТО. Ввод в работу ДТО осуществляется программным ключом \textbf{«ДТО»}. Коэффициент возврата ДТО составляет не менее 0,9.
\end{enumerate}

% ============================================================
% МЕДЛЕННОДЕЙСТВУЮЩИЙ ДИФФЕРЕНЦИАЛЬНЫЙ ОРГАН
% ============================================================
\item Медленнодействующий дифференциальный орган
\begin{enumerate}
\item МДО предназначен для выявления внутренних повреждений, а также для распознавания перехода из внешнего КЗ во внутреннее, в том числе в условиях насыщения ТТ. Ввод в работу МДО осуществляется программным ключом \textbf{«МДО»}.
\item Принцип действия органа основан на контроле формы абсолютного (выпрямленного) мгновенного дифференциального тока. При внутреннем КЗ мгновенный дифференциальный ток изменяется два раза за период промышленной частоты, в то время как при внешнем КЗ с насыщением ТТ – один раз за период. МДО обеспечивает требование по быстродействию при переходе внешнего КЗ во внутреннее.
\item Выпрямленный мгновенный дифференциальный ток рассчитывается как модуль суммы мгновенных значений токов присоединений, входящих в зону действия защиты:
  \begin{align}
    i_\textit{a д}=\mid \sum_{n=1}^N i_{a, n}^{'} \mid,\quad
    i_\textit{b д}=\mid \sum_{n=1}^N i_{b, n}^{'} \mid,\quad
    i_\textit{c д}=\mid \sum_{n=1}^N i_{c, n}^{'} \mid,
  \end{align}
  \begin{eqexpl}[15mm]
    \item{$ i_\textit{a д}, i_\textit{b д}, i_\textit{c д} $} мгновенные значения дифференциальных токов фаз A, B и C, о.е.;
    \item {$ N $} количество присоединений, зафиксированных за соответствующим дифференциальным органом;
    \item {$ i_{a, n}^{'}, i_{b, n}^{'}, i_{c, n}^{'} $} мгновенные значения выровненных фазных токов n-го присоединения, о.е.
  \end{eqexpl}
\end{enumerate}

% ============================================================
% БЫСТРОДЕЙСТВУЮЩИЙ ДИФФЕРЕНЦИАЛЬНЫЙ ОРГАН
% ============================================================
\item Быстродействующий дифференциальный орган
\begin{enumerate}
\item БДО предназначен для быстрого выявления внутренних повреждений до наступления насыщения ТТ и обеспечивает время срабатывания ДЗЛ при внутренних КЗ без выдержки времени не более 25 мс. Ввод в работу данной функции осуществляется программным ключом \textbf{«БДО»}.
\item Принцип действия органа основан на различии скорости нарастания дифференциального и тормозного токов при внутренних и внешних повреждениях. При внутреннем КЗ обеспечивается равный рост дифференциального и тормозного токов в промежутке правильной трансформации ТТ. При внешнем КЗ нарастание тормозного тока происходит раньше, тогда как рост дифференциального тока обусловлен насыщением ТТ.
\item В алгоритме БДО реализована логика выявления одновременного нарастания фронтов действующих значений дифференциального и тормозного токов. Сигнал срабатывания БДО формируется при превышении дифференциального тока уставки срабатывания ДТЗт (\textbf{«Iд нач.»} или \textbf{«Iд нач. груб.»}) и нарастании фронтов дифференциального и тормозного токов в пределах заданного времени.
\end{enumerate}

% ============================================================
% ОРГАН БЛОКИРОВКИ ПРИ ВНЕШНИХ КЗ
% ============================================================
\item Орган блокировки при внешних коротких замыканиях

\begin{enumerate}

\item БВКЗ предназначена для выявления внешнего КЗ, в том числе в условиях насыщения ТТ, и действует на блокировку ДТЗт и ДТО, обеспечивая время достоверного измерения при внешних КЗ на уровне 6 мс. Ввод в работу данной функции осуществляется программым ключом \textbf{«БВКЗ»}.

\item БВКЗ является измерительным органом с зависимой характеристикой дифференциального и тормозного токов. Зоной срабатывания БВКЗ является область под зависимой характеристикой. Данная зона БВКЗ выбрана, чтобы фиксировать внешнее КЗ до наступления насыщения ТТ и появления значительного дифференциального тока, вызванного выходом ТТ плеча ДЗЛ из класса точности.

\item Характеристика срабатывания БВКЗ представлена на рисунке \ref{pict_dzl:BVKZ}.

\dzlfig[width=0.65\textwidth,height=0.65\textheight,keepaspectratio]{\fbpath/02.01 ДЗЛ/561.0651 ДЗЛ/img/BVKZ.pdf}{Характеристика срабатывания БВКЗ}{pict_dzl:BVKZ}

\item Срабатывание БВКЗ происходит при превышении тормозным током 2 о.е. Коэффициент наклона зоны срабатывания БВКЗ выбирается таким образом, чтобы исключить блокировку ДЗЛ при внутренних КЗ, а именно равным коэффициенту наклона характеристики первого наклонного участка ДТЗт.

\item Коэффициент наклона характеристики возврата БВКЗ выбирается таким образом, чтобы исключить возврат блокировки при погрешности ТТ до 70\% в зоне близкой к уровню дифференциального тока при внутреннем КЗ и составляет 1,4 о.е. Коэффициент возврата характеристики БВКЗ для вертикального участка составляет 0,9, при этом предусмотрен автоматический возврат через 300 мс при отсутствии условий пуска.
\end{enumerate}

% ============================================================
% КОНТРОЛЬ ЦЕПЕЙ ТОКА
% ============================================================
\item Контроль цепей тока
\begin{enumerate}
\item В устройстве реализован контроль исправности токовых цепей для исключения ложного срабатывания ДЗЛ вследствие нарушения баланса токов в дифференциальной цепи при обрыве, замыканий и т.д.
\item Ввод в работу контроля исправности токовых цепей осуществляется программным ключом \textbf{«Реж. КЦТ»}. КЦТ срабатывает при превышении дифференциальным током уставки \textbf{«Iср КЦТ»} через выдержку времени \textbf{«Тср КЦТ»}.
  \dzlfig[width=0.9\textwidth,height=0.9\textheight,keepaspectratio]{\fbpath/02.01 ДЗЛ/561.0651 ДЗЛ/img/KCT.pdf}{Контроль цепей тока}{pict_dzl:KCT}
\item Параметры для настройки КЦТ представлены в таблице \ref{app_set:dzl}.
\end{enumerate}

% ============================================================
% БЛОКИРОВКА ДЗЛ ОТ ЦКС
% ============================================================
\item Блокировка ДЗЛ от ЦКС
\begin{enumerate}
\item ФБ Блокировка ДЗЛ от ЦКС предназначен для формирования сигнала блокировки ДЗЛ при неисправности ЦКС или при приеме внешних сигналов блокировки от удаленных терминалов (УТ) ДЗЛ по ЦКС1 и ЦКС2.
\item При реализации ДЗЛ двухконцевой линии ДЗЛ блокируются при одновременной неисправности обоих ЦКС, для трехконцевой ДЗЛ – при неисправности одного из ЦКС. Конфигурация защищаемой линии задается программным переключателем \textbf{«Топология ЛЭП»}. При этом для трехконцевой линии действие на отключение в этом случае формируется при приеме сигналов срабатывания ДЗЛ от п/к, сохранившего связь со с удаленными.
  \dzlfig[width=0.9\textwidth,height=0.9\textheight,keepaspectratio]{\fbpath/02.01 ДЗЛ/561.0651 ДЗЛ/img/Blk_DZL_CKS.pdf}{Блокировка ДЗЛ от ЦКС}{pict_dzl:Blk_DZL_CKS}
\item Параметры для настройки блокировки ДЗЛ от ЦКС представлены в таблице \ref{app_set:dzl}.
\end{enumerate}

% ============================================================
% ДИФФЕРЕНЦИАЛЬНАЯ ТОКОВАЯ ЗАЩИТА С ТОРМОЖЕНИЕМ (логика)
% ============================================================
\item Дифференциальная токовая защита с торможением (логика)
\begin{enumerate}
\item Функциональный блок ДТЗт предназначен для формирования пофазных и общего сигналов пуска, срабатывания и блокировки ДТЗт.
\item Функциональная схема логики ДТЗт представлена на рисунке \ref{pict_dzl:DTZt}.
  \dzlfig[width=0.9\textwidth,height=0.9\textheight,keepaspectratio]{\fbpath/02.01 ДЗЛ/561.0651 ДЗЛ/img/DTZt.pdf}{Функциональная схема ДТЗт}{pict_dzl:DTZt}
\item Пуск ДТЗт осуществляется при срабатывании ИО ДТЗт с контролем срабатывания специальных ИО (МДО, БДО). Срабатывание ДТЗт осуществляется при сохранении условий пуска в течении выдержки времени \textbf{«Тср ДТЗт»}.
\item Блокировка ДТЗт осуществляется при срабатывании БВКЗ. Предусмотрена возможность блокировки ДТЗт в случае выявления неисправности в токовых цепях при введенном программном ключе \textbf{«Контр. ДТЗт от КЦТ»}, в случае выявления неисправности ЦКС или при блокировке ДЗЛ в УТ при введенном программном ключе \textbf{«Контр. ДТЗт от ЦКС»}, при протекании сквозных токов внешнего КЗ через ошиновку при введенном программном ключе \textbf{«БСТО»}.
\item Предусмотрена возможность контроля ОПТО при введенном программном ключе \textbf{«Контр. ДТЗт от ОПТО»}. ОПТО обеспечивает селективную работу ДТЗт на линиях с ответвлениями, разрешая работу ДТЗт только при КЗ на защищаемой линии.
\item Параметры для настройки ДТЗт представлены в таблице \ref{app_set:dzl}.
\end{enumerate}

% ============================================================
% ДИФФЕРЕНЦИАЛЬНАЯ ТОКОВАЯ ОТСЕЧКА (логика)
% ============================================================
\item Дифференциальная токовая отсечка (логика)
\begin{enumerate}
\item Функциональный блок ДТО предназначен для формирования пофазных и общего сигналов пуска, срабатывания и блокировки ДТО.
\item Пуск ДТО осуществляется при срабатывании ИО ДТО. Срабатывание ДТО осуществляется при сохранении условий пуска в течении выдержки времени \textbf{«Тср ДТО»}.
\item Блокировка ДТО осуществляется при срабатывании БВКЗ. Предусмотрена возможность блокировки ДТО в случае выявления неисправности ЦКС или при блокировке ДЗЛ в УТ при введенном программном ключе \textbf{«Контр. ДТЗт от ЦКС»}, при протекании сквозных токов внешнего КЗ через ошиновку при введенном программном ключе \textbf{«БСТО»}.
\item Функциональная схема логики ДТО представлена на рисунке \ref{pict_dzl:DTO}.
  \dzlfig[width=0.9\textwidth,height=0.9\textheight,keepaspectratio]{\fbpath/02.01 ДЗЛ/561.0651 ДЗЛ/img/DTO.pdf}{Функциональная схема ДТО}{pict_dzl:DTO}
\item Параметры для настройки ДТО представлены в таблице \ref{app_set:dzl}.
\end{enumerate}

% ============================================================
% ЛОГИКА СРАБАТЫВАНИЯ ДЗЛ
% ============================================================
\item Логика срабатывания ДЗЛ
\begin{enumerate}
\item Логика срабатывания ДЗЛ предназначена для формирования сигналов срабатывания ДЗЛ, а также сигналов блокировки ДЗЛ удаленных терминалов.
\item Действие на отключение от ДЗЛ формируется при срабатывании ДТЗт или ДТО, а также при приеме внешних сигналов срабатывания ДЗЛ от УТ ДЗЛ по ЦКС1 и ЦКС2.
\item Блокировка ДЗЛ в удаленных терминалах осуществляется при оперативном выводе функции ДЗЛ или терминала, при вводе режима тестирования локального п/к или при срабатывании БСТО при введенном программном ключе \textbf{«БСТО ДЗЛ»}. Предусмотрена возможность блокировки ДЗЛ в удаленных терминалах при переводе ДЗЛ на сигнал при введенном программном ключе \textbf{«Блок. ДЗЛ УТ при переводе на сигнал»}.
  \dzlfig[width=0.9\textwidth,height=0.9\textheight,keepaspectratio]{\fbpath/02.01 ДЗЛ/561.0651 ДЗЛ/img/Log_DZL.pdf}{Логика срабатывания ДЗЛ}{pict_dzl:Log_DZL}
\item Параметры для настройки логики срабатывания ДЗЛ представлены в таблице \ref{app_set:dzl}.
\end{enumerate}

% ============================================================
% ЦИФРОВОЙ КАНАЛ СВЯЗИ ДЗЛ (обмен данными)
% ============================================================
\item Цифровой канал связи ДЗЛ (обмен данными)
\begin{enumerate}
\item Обмен данными между п/к осуществляется по цифровым каналам связи. Каждый п/к оснащён двумя независимыми портами передачи данных, предназначенными для подключения к другим п/к защиты. В зависимости от конфигурации защищаемой линии (двух- или трёхконцевой), структура связи реализуется по-разному:
  \begin{itemize}
    \item для двухконцевой линии один порт каждого п/к используется для подключения по основному каналу связи, а второй — по резервному. Связь между п/к в данном случае осуществляется по схеме «ведущий-ведущий». При этом обеспечивается автоматическое переключение между основным и резервным портами связи при потере связи по основному каналу;
    \item для трёхконцевой линии п/к соединяются в кольцевую топологию, при которой каждый из них обменивается данными с двумя другими. При нормальной работе каналов связи каждый из трёх п/к получает цифровую информацию о токах c двух других концов линии и выполняет расчет дифференциальной и тормозной составляющих на основе собственного тока и токов, полученных от смежных п/к. При выполнении условий срабатывания каждый п/к действует на отключение своего выключателя. В случае разрыва связи между двумя терминалами (при сохранении связи с третьим), расчет дифференциального и тормозного токов выполняется п/к, сохранившим связь с обоими удаленными. При выполнении условий срабатывания действие удаленных п/к на отключение производится посредством передачи на них сигналов отключения \textit{«Ср. ДЗЛ ф.А УТ по ЦКС»}, \textit{«Ср. ДЗЛ ф.В УТ по ЦКС»}, \textit{«Ср. ДЗЛ ф.С УТ по ЦКС»}.
  \end{itemize}
\item Осуществляется автоматический контроль исправности канала связи между п/к. При обнаружении неисправностей формируются сигналы \textbf{«Неиспр. ЦКС1»} и \textbf{«Неиспр. ЦКС2»}.
\end{enumerate}

% ============================================================
% ОСОБЕННОСТИ ОРГАНИЗАЦИИ ЦКС ДЗЛ
% ============================================================
\item Особенности организации цифрового канала связи ДЗЛ

\begin{enumerate}

\item При организации связи между п/к ДЗЛ по выделенному оптическому каналу связи предпочтительным является использование пары волокон, когда одна жила используется для передачи данных в одном направлении канала, а вторая жила -- в противоположном. При организации ВОЛС большой протяженности рекомендуется применение оптических усилителей для обеспечения системного запаса по затуханию, при этом максимальная длина ВОЛС определяется типом применяемого оптического усилителя.

\item При организации связи между п/к ДЗЛ с использованием оборудования ЦСПИ по мультиплексированному каналу связи подключение п/к ДЗЛ к мультиплексору должно осуществляться по оптическому интерфейсу C37.94. При наличии оптического интерфейса C37.94 в мультиплексоре возможен вариант подключения п/к ДЗЛ напрямую к мультиплексору по ВОЛС на расстоянии до 2 км. В случае использования электрических интерфейсов E1, X.21 или G.703.1 в мультиплексоре рекомендуется применение специального конвертора между п/к ДЗЛ и мультиплексором для преобразования оптических сигналов в электрические.

\item Варианты организации одного канала связи между п/к ДЗЛ приведены на рисунках ниже. Работа второго канала может организовываться аналогично первому, либо возможен вариант работы одного канала по выделенному оптическому каналу связи, а другого -- через ЦСПИ.

\item При выполнении защит линии с использованием трех п/к ДЗЛ допускается использовать один КС для обеспечения взаимодействия двух п/к ДЗЛ и рекомендуется подключение трех п/к ДЗЛ в кольцо для повышения надежности.

\dzlfig[width=0.9\textwidth,height=0.9\textheight,keepaspectratio]{\fbpath/02.01 ДЗЛ/561.0651 ДЗЛ/img/Cnl_1.pdf}{Организация связи между двумя п/к ДЗЛ по выделенному оптическому каналу связи с использованием двух жил оптоволокна}{pict_dzl:Cnl_1}

\dzlfig[width=0.9\textwidth,height=0.9\textheight,keepaspectratio]{\fbpath/02.01 ДЗЛ/561.0651 ДЗЛ/img/Cnl_2.pdf}{Организация связи между двумя п/к ДЗЛ по выделенному оптическому каналу связи через оптические усилители}{pict_dzl:Cnl_2}

\dzlfig[width=0.9\textwidth,height=0.9\textheight,keepaspectratio]{\fbpath/02.01 ДЗЛ/561.0651 ДЗЛ/img/Cnl_3.pdf}{Организация связи между двумя п/к ДЗЛ через ЦСПИ}{pict_dzl:Cnl_3}

\dzlfig[width=0.9\textwidth,height=0.9\textheight,keepaspectratio]{\fbpath/02.01 ДЗЛ/561.0651 ДЗЛ/img/Cnl_4.pdf}{Организация связи между двумя п/к ДЗЛ через ЦСПИ с использованием преобразователей интерфейса}{pict_dzl:Cnl_4}

\dzlfig[width=0.9\textwidth,height=0.9\textheight,keepaspectratio]{\fbpath/02.01 ДЗЛ/561.0651 ДЗЛ/img/Cnl_5.pdf}{Организация связи между тремя п/к ДЗЛ по выделенному оптическому каналу связи с использованием двух жил оптоволокна}{pict_dzl:Cnl_5}

\dzlfig[width=0.9\textwidth,height=0.9\textheight,keepaspectratio]{\fbpath/02.01 ДЗЛ/561.0651 ДЗЛ/img/Cnl_6.pdf}{Организация связи между тремя п/к ДЗЛ через ЦСПИ}{pict_dzl:Cnl_6}

\item В случае обнаружения неисправности одного из каналов связи с удаленным п/к ДЗЛ защита переключается на другой канал связи. При возникновении КЗ на защищаемой линии и одновременном повреждении одного из подключенных к п/к ДЗЛ каналов связи применяемый алгоритм резервирования каналов связи обеспечивает время срабатывания защиты не более 40 мс.
\end{enumerate}

% ============================================================
% ТРЕБОВАНИЯ К ЦКС ДЗЛ
% ============================================================
\item Требования к цифровому каналу связи ДЗЛ
\begin{enumerate}
\item Для корректной работы ДЗЛ канал связи (с учетом работы оборудования ЦСПИ) должен соответствовать следующим требованиям:
  \begin{itemize}
    \item максимальная допустимая задержка передачи данных в одном направлении не более 30 мс;
    \item максимальная допустимая асимметрия задержки передачи данных (разность времени передачи данных в прямом и обратном направлениях) не более 0,5 мс;
    \item одновременное в обоих направлениях (на прием и на передачу) переключение КС на резервный маршрут при потере основного маршрута или обеспечение статического маршрута без автоматического переключения при передаче данных через мультиплексированные КС.
  \end{itemize}
\end{enumerate}

% ============================================================
% ПАРАМЕТРЫ НАСТРОЙКИ ЦКС ДЗЛ
% ============================================================
\item Параметры настройки цифрового канала связи ДЗЛ
\begin{enumerate}
\item Параметр \textbf{«Синх. КС»} задает источник синхронизации передаваемых данных в КС и должен быть определен исходя из его типа (выделенный или мультиплексированный). Для синхронизации данных может быть использован синхросигнал от внутреннего тактового генератора (Внутренняя) либо синхронизирующий сигнал, выделяемый из принимаемых данных локального п/к ДЗЛ (Внешняя). В случае применения мультиплексированных КС параметр \textbf{«Синх. КС»} должен устанавливаться в значение \textit{Внешняя} независимо от способа подключения к мультиплексору. Обязательным требованием к мультиплексированному КС является синхронность по каналу передачи и каналу приема, что должно обеспечиваться соответствующими настройками оборудования связи. В случае применения выделенных КС отсутствует необходимость использовать внешнюю синхронизацию по причине отсутствия дополнительных устройств в канале связи, с которыми необходимо синхронизировать данные. Поэтому параметр \textbf{«Синх. КС»} должен быть установлен в значение \textit{Внутренняя}. При использовании кодирования, предусмотренного стандартом С37.94 на выделенном канале связи возможно использование как внутренней, так и внешней синхронизации. Таким образом один из п/к ДЗЛ может быть источником сигнала синхронизации, а другой п/к ДЗЛ может быть приемником сигнала синхронизации.
\item Параметр \textbf{«Скор»} на выделенных КС устанавливает логическую скорость при использовании внутренней синхронизации. При использовании внешней синхронизации параметр игнорируется. При использовании мультиплексированного КС параметр \textbf{«Скор»} игнорируется и может быть выставлен любым, так как скорость передачи данных выставляется непосредственно в мультиплексоре установкой числа N, которое определяет количество 64-килобитных каналов, используемых для формирования общей скорости передачи. По умолчанию параметр \textbf{«Скор»} равен 64 кБит/с. При наличии двух плеч на присоединение параметр \textbf{«Скор»} должен быть равен не менее 128 кБит/с.
\item Параметр \textbf{«Комп. асим.»} необходим для компенсации постоянной составляющей асимметрии КС. На выделенных КС параметр \textbf{«Комп. асим.»} устанавливается равным 0, так как длина ВОЛС в прямом и обратном направлении, как правило, одинакова. Поэтому выделенные КС принципиально симметричны. Для мультиплексированных КС необходимо задать режим п/к ДЗЛ уставкой \textbf{«Реж. терм.»} При этом необходимо правильно устанавливать режим каждого п/к ДЗЛ. Например, в случае установки в обоих п/к ДЗЛ режима «Ведущий» каждый п/к ДЗЛ использует собственное установленное время для компенсации асимметрии. В случае установки одинакового режима в обоих п/к ДЗЛ компенсация асимметрии будет работать неправильно. П/к ДЗЛ с режимом «Ведущий» использует заданное значение параметра \textbf{«Комп. асим.»} и передает данное значение в удаленный п/к ДЗЛ. П/к ДЗЛ с режимом «Ведомый» игнорирует значение параметра \textbf{«Комп. асим.»} и использует принятое значение от удаленного п/к ДЗЛ с противоположным знаком.
\item Параметры \textbf{«ID Tx»} и \textbf{«ID Rx»} предназначены для исключения ошибок при подключении каналов связи вследствие неправильной коммутации потоков данных. Использование уникальных идентификаторов для каждого КС (для приема и передачи данных) позволяет выявить подобные ошибки и исключить неправильную работу функции ДЗЛ. П/к ДЗЛ проверяет соответствие установленного параметра \textbf{«ID Rx»} и принимаемого идентификатора в данных от удаленного п/к ДЗЛ. В случае их несоответствия формируется сигнал неисправности канала связи.
\end{enumerate}

% ============================================================
% ДИАГНОСТИКА СОСТОЯНИЯ ЦКС
% ============================================================
\item Диагностика состояния цифрового канала связи
\begin{enumerate}
\item Каждый п/к ДЗЛ осуществляет автоматический непрерывный контроль состояния каналов связи. Для каждого канала связи (ЦКС1 и ЦКС2) осуществляется регистрация в энергонезависимой памяти терминала и отображение на ИЧМ терминала диагностических сигналов для контроля исправности, а также проверки правильности задания параметров настройки.
\item В таблице ниже приведен перечень диагностических сигналов ЦКС1, для ЦКС2 перечень аналогичный.
  \begingroup
  \small
  \setlength{\extrarowheight}{0.06cm}
  \begin{longtable}{|>{\centering\arraybackslash}m{4cm}|>{\centering\arraybackslash}m{6.4cm}|>{\centering\arraybackslash}m{2.8cm}|>{\centering\arraybackslash}m{2.8cm}|}
    \caption{Перечень диагностических сигналов ЦКС1\hfill\vspace{-0.5\baselineskip}}\label{desc_dzl:tbl10}\\
    \hline
    \rowcolor{gray!30} Обозначение & Описание & Знач. при исправном состоянии & Знач. при неисправ. состоянии \\
    \hline
    \endfirsthead
    \caption*{\hspace{3pt}\emph{Продолжение таблицы \ref{desc_dzl:tbl10}}\hfill\ }\vspace{-0.5\baselineskip} \\
    \hline
    \rowcolor{gray!30} Обозначение & Описание & Знач. при исправном состоянии & Знач. при неисправ. состоянии \\
    \hline
    \endhead
    \centering\arraybackslash ЦКС1: Приним. ID  & \centering\arraybackslash Принимаемый идентификатор устройства ЦКС1 & \centering\arraybackslash =ID Rx & \centering\arraybackslash ≠ ID Rx \\
    \hline
    \centering\arraybackslash ЦКС1: t\_ПЕР  & \centering\arraybackslash Текущее значение времени задержки передачи данных по ЦКС1, мкс & \centering\arraybackslash 0-30000 & \centering\arraybackslash >30000 \\
    \hline
    \centering\arraybackslash ЦКС1: t\_АС  & \centering\arraybackslash Измеренное значение времени асимметрии задержки передачи данных по ЦКС1 (при наличии внешней синхронизации), мкс & \centering\arraybackslash 0 & \centering\arraybackslash >|±500| \\
    \hline
    \centering\arraybackslash ЦКС1: Сбой синхр.  & \centering\arraybackslash Неисправность внешнего источника синхронизации  & \centering\arraybackslash 0 & \centering\arraybackslash 1 \\
    \hline
    \centering\arraybackslash ЦКС1: Уд. сбой синхр.  & \centering\arraybackslash Неисправность внешнего источника синхронизации на удаленном устройстве & \centering\arraybackslash 0 & \centering\arraybackslash 1 \\
    \hline
    \centering\arraybackslash ЦКС1: Отсут. приема  & \centering\arraybackslash Отсутствие приема данных по ЦКС & \centering\arraybackslash 0 & \centering\arraybackslash 1 \\
    \hline
    \centering\arraybackslash ЦКС1: Уд. отсут. приема  & \centering\arraybackslash Отсутствие приема данных по ЦКС1 на удаленном устройстве & \centering\arraybackslash 0 & \centering\arraybackslash 1 \\
    \hline
    \centering\arraybackslash ЦКС1: Неверный ID  & \centering\arraybackslash Прием данных с неверным адресом источника по ЦКС1 & \centering\arraybackslash 0 & \centering\arraybackslash 1 \\
    \hline
    \centering\arraybackslash ЦКС1: Сбой данных  & \centering\arraybackslash Неверная контрольная сумма принятых данных по ЦКС1 & \centering\arraybackslash 0 & \centering\arraybackslash 1 \\
    \hline
    \centering\arraybackslash ЦКС1: Сбой передачи  & \centering\arraybackslash Недопустимая задержка передачи данных по ЦКС1 & \centering\arraybackslash 0 & \centering\arraybackslash 1 \\
    \hline
    \centering\arraybackslash ЦКС1: Сч. ошибок  & \centering\arraybackslash Текущее количество ошибок в ЦКС1 & \centering\arraybackslash 0--10000, не изменяется & \centering\arraybackslash Увелич. на 1 \\
    \hline
    \centering\arraybackslash ЦКС1: Неисправность КС  & \centering\arraybackslash Неисправное состояние ЦКС1 & \centering\arraybackslash 0 & \centering\arraybackslash 1 \\
    \hline
  \end{longtable}
  \endgroup
\item Параметр \textit{Приним. ID} показывает идентификатор устройства КС, получаемый в данных от удаленного п/к ДЗЛ. При исправном состоянии указанный параметр соответствует заданному значению \textbf{«ID Rx»}.
\item Параметр \textit{t\_ПЕР} показывает измеряемое время задержки передачи данных в одном направлении, определяемое скоростью передачи данных, длиной канала связи, а также для мультиплексированных КС количеством мультиплексоров и способом переприема данных в них. При исправном состоянии указанный параметр не превышает 30 мс.
\item Параметр \textit{t\_АС} показывает измеряемое время асимметрии задержки передачи данных в прямом и обратном направлении. При отсутствии сигнала 1PPS хотя бы в одном из п/к ДЗЛ указанный параметр должен устанавливаться равным нулю, что не является признаком симметричности КС. Измерение может осуществляться только при наличии сигналов 1PPS от устройств синхронизации GPS/ГЛОНАСС в обоих п/к ДЗЛ. При исправном состоянии указанный параметр не превышает 0,5 мс. Значения отображаются с учетом выставленной уставки компенсации асимметрии \textbf{«Комп. асим.»}.
\item Сигнал неисправности \textit{«Отсут. приема»} формируется при отсутствии данных на приемнике порта локального п/к ДЗЛ от удаленного устройства (п/к ДЗЛ или оборудования ЦСПИ).
\item Сигнал неисправности \textit{«Уд. отсут. приема»} формируется при отсутствии данных на приемнике порта удаленного устройства (п/к ДЗЛ или оборудования ЦСПИ) от локального п/к ДЗЛ.
\item Сигнал неисправности \textit{«Неверный ID»} формируется при несоответствии принимаемого идентификатора КС удаленного п/к ДЗЛ заданному значению \textbf{«ID Rx»}.
\item Сигнал неисправности \textit{«Сбой данных»} формируется при неверной контрольной сумме принятых данных, что означает превышение вероятности битовых ошибок в канале связи над требуемым значением для штатной работы п/к ДЗЛ (для канала связи через ЦСПИ - $ 10^{-3}) $.
\item Сигнал неисправности \textit{«Сбой передачи»} формируется при превышении времени задержки передачи данных максимально допустимого значения, равного 30 мс.
\item Сигнал \textit{«Сч. Ошибок»} не является признаком неисправности и показывает общее количество ошибок в КС, возникших до текущего момента времени. При фиксации одной из ошибок в КС значение счетчика ошибок увеличивается на единицу. Значение счетчика ошибок не сбрасывается при перерыве питания терминала.
\item Общий сигнал неисправности КС формируется при фиксации одной из неисправностей в КС (\textit{«Отсут. приема», «Уд. отсут. Приема», «Неверный ID», «Сбой данных», «Сбой передачи», «Сбой синхр.», «Уд. сбой синхр.»}) в течение фиксированной выдержки времени 50 мс или при наличии соответствующего входного внешнего сигнала неисправности КС.
\end{enumerate}

\end{enumerate}